\chapter*{TÓM TẮT KHOÁ LUẬN}
\addcontentsline{toc}{section}{{\bf TÓM TẮT KHOÁ LUẬN}\rm}

Sự xuất hiện và phát triển không ngừng của các giải pháp bảo mật đã đóng góp một phần không nhỏ trong việc bảo đảm sự an toàn của chúng ta trước những cuộc tấn công từ các người dùng, tổ chức không thiện chí trên không gian mạng. Tuy nhiên, sự phát triển của các giải pháp bảo mật cũng kéo theo sự tiến hóa của các cuộc tấn công. Trong đó, tấn công APT là một trong những hình thức tấn công mới, cực kỳ tinh vi mà chúng ta phải đối mặt. Trong các cuộc tấn công dạng này, các kẻ tấn công thường sử dụng những kỹ thuật vô cùng tân tiến, hiện đại để xâm nhập vào hệ thống của mục tiêu và lẩn tránh khỏi các giải phải pháp bảo mật hiện hành. Đồng thời, các kẻ tấn công cũng sẽ giới hạn hoạt động của mình kết hợp với việc sử dụng các kỹ thuật lẩn tránh tinh vi nhằm tồn tại trong hệ thống mục tiêu lâu nhất có thể. Chính vì thế, các cuộc tấn công loại này thường rất khó phát hiện và gây ra những thiệt hại khổng lồ cho các tổ chức bị nhắm đến.
\bigbreak\noindent
Trong khóa luận này, nhằm mục tiêu đề ra một giải pháp hiệu quả trong việc phòng chống và phát hiện tấn công APT, chúng tôi đã tiến hành xây dựng một hệ thống săn tìm mối đe dọa ứng dụng AI trong việc tự động hóa quá trình phát hiện các hoạt động bất thường có thể dẫn đến một cuộc tấn cuộc tấn công APT. Hệ thống được đề xuất sử dụng các mạng nơ-ron khác nhau như CNN, GRU, GNN, ... trong các tác vụ chuyên biệt nhằm mang đến hiệu quả phát hiện bất thường cao nhất. Bên cạnh đó, phương pháp học liên kết được ứng dụng nhằm nâng cao khả năng chia sẻ thông tin giữa các tổ chức mà không cần phải quan ngại về vấn đề bảo mật và tính riêng tư của dữ liệu. Hơn thế nữa, các học máy khả giải và một phương pháp xác định độ tin cậy của dự đoán từ mô hình AI được sử dụng để hỗ trợ cho các nhà phân tích bảo mật trong việc hiểu và xác định tính đúng đắn của các quyết định từ các mô hình học sâu. Cuối cùng, hệ thống được thiết kế với sự tham gia của các giải pháp bảo mật hiện đại có ý nghĩa quan trọng trong săn tìm mối đe dọa như SIEM, CTI, ... và được triển khai trong môi trường mạng SDN nhằm tạo ra một giải pháp toàn diện, hiện đại và có khả năng ứng dụng trong thực tế. Hệ thống được đề xuất được chúng tôi đặt tên là \textbf{XFedHunter}.
\bigbreak\noindent
Tính hiệu quả của các mô hình học sâu được huấn luyện thông qua phương pháp học liên kết và các giải pháp học máy khả giải được đề xuất được chúng tôi kiểm nghiệm thông qua hai tập dữ liệu là NF-ToN-IoT và DARPA TCE3. Ngoài ra, một hệ thống săn tìm tự động trong mạng SDN cũng được chúng tôi triển khai nhằm chứng minh tính khả thi và thực tế của hệ thống đề xuất.
